\documentclass[10pt,a4paper]{article}
\usepackage[T1]{fontenc}
\usepackage[utf8]{inputenc}
\usepackage{tgtermes,tgheros,microtype}
\usepackage[a4paper,margin=18mm,headheight=19pt]{geometry}
\usepackage{xcolor,graphicx,booktabs,tabularx,array,enumitem}
\usepackage[most]{tcolorbox}
\usepackage{fancyhdr,lastpage,titlesec}
\usepackage[hidelinks]{hyperref}
\usepackage{xurl,caption,placeins,textcomp,amsmath,longtable}

\definecolor{Navy}{HTML}{173753}
\definecolor{Teal}{HTML}{008F8C}
\definecolor{PaleTeal}{HTML}{E5F4F3}
\definecolor{Orange}{HTML}{F08A24}
\definecolor{PaleOrange}{HTML}{FFF1E2}
\definecolor{Red}{HTML}{BB3E35}
\definecolor{PaleRed}{HTML}{F9E9E7}
\definecolor{Slate}{HTML}{5B7083}
\definecolor{Light}{HTML}{F3F6F8}
\definecolor{Line}{HTML}{D7E0E7}
\definecolor{Green}{HTML}{2E8B57}
\definecolor{PaleGreen}{HTML}{EAF5EE}

\newcommand{\sans}{\fontfamily{qhv}\selectfont}
\newcommand{\reportdate}{31 August 2026}
\newcommand{\reportversion}{Version 9}
\newcolumntype{Y}{>{\raggedright\arraybackslash}X}
\newcolumntype{R}{>{\raggedleft\arraybackslash}X}
\newcolumntype{C}{>{\centering\arraybackslash}X}

\hypersetup{
  pdftitle={HeiQ XMembrane Final Result and Decision Report -- Version 9},
  pdfauthor={HeiQ XMembrane project team},
  pdfsubject={Final integrated ERV+EO model, experiment and delivery basis},
  pdfkeywords={XMembrane, EO, ERV, integrated membrane, O30, break-even, 9000 m2, vapour cell}
}

\pagestyle{fancy}
\fancyhf{}
\fancyhead[L]{\sans\footnotesize\color{Slate}HEIQ MATERIALS AG · XMEMBRANE}
\fancyhead[R]{\sans\footnotesize\color{Slate}FINAL RESULT REPORT · V9}
\fancyfoot[L]{\sans\scriptsize\color{Slate}Conditional modelling and experiment-planning evidence · no product validation}
\fancyfoot[R]{\sans\scriptsize\color{Slate}Page \thepage\ of \pageref*{LastPage}}
\renewcommand{\headrulewidth}{0.45pt}
\renewcommand{\headrule}{\hbox to\headwidth{\color{Navy}\leaders\hrule height \headrulewidth\hfill}}
\renewcommand{\footrulewidth}{0pt}

\titleformat{\section}{\sans\bfseries\color{Navy}\fontsize{19}{22}\selectfont}{\thesection}{0.65em}{}
\titleformat{\subsection}{\sans\bfseries\color{Navy}\fontsize{13}{16}\selectfont}{\thesubsection}{0.55em}{}
\titlespacing*{\section}{0pt}{7mm}{3mm}
\titlespacing*{\subsection}{0pt}{5mm}{2mm}
\setlength{\parindent}{0pt}
\setlength{\parskip}{1.6mm}
\setlist[itemize]{leftmargin=5mm,itemsep=1mm,topsep=1mm}
\captionsetup{font=small,labelfont={bf,color=Navy}}
\renewcommand{\arraystretch}{1.18}

\newtcolorbox{mainbox}{enhanced,breakable,colback=PaleTeal,colframe=Teal,boxrule=0.7pt,arc=2mm,left=4mm,right=4mm,top=3mm,bottom=3mm}
\newtcolorbox{warningbox}{enhanced,breakable,colback=PaleOrange,colframe=Orange,boxrule=0.7pt,arc=2mm,left=4mm,right=4mm,top=3mm,bottom=3mm}
\newtcolorbox{stopbox}{enhanced,breakable,colback=PaleRed,colframe=Red,boxrule=0.7pt,arc=2mm,left=4mm,right=4mm,top=3mm,bottom=3mm}
\newtcolorbox{gobox}{enhanced,breakable,colback=PaleGreen,colframe=Green,boxrule=0.7pt,arc=2mm,left=4mm,right=4mm,top=3mm,bottom=3mm}
\newtcolorbox{plainbox}{enhanced,breakable,colback=Light,colframe=Line,boxrule=0.5pt,arc=2mm,left=4mm,right=4mm,top=3mm,bottom=3mm}
\newcommand{\eyebrow}[1]{{\sans\bfseries\color{Teal}\MakeUppercase{#1}}\par\vspace{1.5mm}}
\newcommand{\metric}[3]{\begin{minipage}[t]{0.30\textwidth}{\sans\bfseries\fontsize{18}{20}\selectfont\color{Navy}#1}\par{\sans\bfseries\small #2}\par{\scriptsize\color{Slate}#3}\end{minipage}}

\begin{document}
\thispagestyle{empty}
\vspace*{7mm}
\begin{tcolorbox}[enhanced,colback=Navy,colframe=Navy,boxrule=0pt,arc=0mm,left=12mm,right=12mm,top=12mm,bottom=12mm]
{\sans\bfseries\large\color{white}HEIQ XMEMBRANE PROJECT TEAM}\par
\vspace{8mm}
{\sans\bfseries\fontsize{28}{32}\selectfont\color{white}Final result and decision\\[-1mm]report}\par
\vspace{2mm}
{\sans\fontsize{18}{21}\selectfont\color{white}EO integrated only in the ERV membrane/core}\par
\vspace{7mm}
{\sans\large\color{white}\reportdate\quad · \quad\reportversion}
\end{tcolorbox}

\vspace{10mm}
{\sans\bfseries\Large\color{Navy}Controlled final project basis}\par
\vspace{4mm}
This report supersedes Version~8 and all earlier architecture comparisons for
current decisions. The proposed EO function is located \textbf{inside the ERV
membrane/core}; no downstream or post-ERV EO module is part of the selected
architecture. Older calculations are retained only as provenance.

\begin{mainbox}
\textbf{Final model-phase conclusion.} The matched model identifies a
conditional incremental energy window for EO integrated in the ERV. O30 is
reachable below the 9,000-m\textsuperscript{2} area gate at 5,899~m\textsuperscript{2},
with 418.400~kW known power and a 125.340-kW (23.05\%) window versus the same
passive ERV+AC system with EO off. This is strong enough to justify material
and coupon testing, but not a validated net-system saving.
\end{mainbox}

\vfill
\begin{tabularx}{\textwidth}{@{}lY@{}}
\textbf{Fixed service} & O80 inlet at 32~$^\circ$C/80\% RH to 20~$^\circ$C/50\% RH supply \\
\textbf{Primary comparator} & Passive ERV+AC = 543.741~kW; AC only = 1,166.400~kW \\
\textbf{Selected integrated point} & O30, 1.00~V, 5,899~m\textsuperscript{2}, 418.400~kW known subtotal \\
\textbf{Calculated EO break-even} & 51.955~Wh/kg all-in active-core energy \\
\textbf{Release status} & Controlled V9; source, selected results, manifest and SHA256 checksums packaged \\
\textbf{Claim status} & Conditional model evidence; no validated membrane, module or commercial saving \\
\end{tabularx}

\clearpage
\section*{Executive summary}
\addcontentsline{toc}{section}{Executive summary}

The largest established model reduction comes from the passive ERV: the
equal-service subtotal falls from 1,166.400~kW for AC alone to 543.741~kW for
passive ERV+AC. EO must therefore be judged incrementally against the same ERV
with EO off. Crediting the entire AC-to-active-system difference to EO would
overstate its contribution.

Passive ERV already reaches 13.215~g/kg, below the O50 coordinate of
14.952~g/kg. O50 therefore has zero EO duty in the matched model. At O30
(8.886~g/kg), the selected area-constrained integrated ERV+EO point uses
5,899~m\textsuperscript{2}. It combines 316.910~kW of remaining AC with
101.490~kW of known active-core power, for 418.400~kW total. The incremental
window is 125.340~kW versus passive ERV+AC and 747.999~kW versus AC alone; only
the former is the EO-specific comparison.

The O30 break-even is calculated, not assumed:
\[
e_{\mathrm{BE}}=
\frac{P_{\mathrm{passive\ ERV+AC}}-P_{\mathrm{remaining\ AC}}}
{\dot m_{\mathrm{water,integrated\ core}}}
=51.955~\mathrm{Wh/kg}.
\]
The known active-core SEC is 23.246~Wh/kg, leaving 28.709~Wh/kg or 125.340~kW
for omitted receiver, heat-rejection and hardware loads before the incremental
advantage disappears. The historical 180-Wh/kg shorthand is not the
architecture-specific break-even.

Full humidity removal is not available below the current area gate. The
minimum valid screened full-duty point requires approximately
17,253~m\textsuperscript{2} and 113.7~kW known total power. The lowest-power
stored full-duty mechanism coordinate uses 34,425~m\textsuperscript{2} and
107.856~kW. Both are outside the accepted area envelope. The practical
opportunity is therefore partial active enhancement of the ERV, with O30 the
best current under-gate coordinate, rather than forcing EO to complete all
dehumidification.

\begin{center}
\metric{5,899 m\textsuperscript{2}}{Selected O30 area}{Inside the 9,000-m\textsuperscript{2} gate}\hfill
\metric{125.3 kW}{Incremental window}{Versus the same ERV with EO off}\hfill
\metric{52.0 Wh/kg}{Calculated break-even}{Known active core: 23.2 Wh/kg}
\end{center}

\clearpage
\tableofcontents
\clearpage

\section{Architecture and comparison discipline}
\eyebrow{One inlet, one endpoint, one matched passive control}

The active and passive cases share the same outdoor inlet, supply endpoint,
airflow and sensible-effectiveness basis. EO acts only across the active
membrane inside the ERV core. The downstream coil removes the remaining
sensible and latent load; at the full humidity endpoint it is sensible only.
O50 and O30 are intermediate humidity-ratio coordinates, not alternative
weather states and not supply-air definitions.

\begin{figure}[htbp]
  \centering
  \includegraphics[width=0.98\textwidth]{../../outputs/integrated_erv_eo_opportunity_plots_v2/integrated_power_ladder.pdf}
  \caption{Matched equal-service architectures. EO exists only in the ERV membrane/core. The O30 selection is below 9,000~m\textsuperscript{2}; the full-duty point is an unconstrained mechanism coordinate.}
\end{figure}
\FloatBarrier

\begin{warningbox}
\textbf{Comparator rule.} Use AC+ERV versus AC to quantify passive recovery.
Use AC+ERV+EO versus the same AC+ERV case to quantify incremental EO benefit.
The AC-only comparison may be reported separately, but must not be labelled as
an EO-only saving.
\end{warningbox}

\begin{figure}[htbp]
  \centering
  \includegraphics[width=0.97\textwidth]{../../outputs/integrated_erv_eo_opportunity_plots_v2/integrated_pairwise_reductions.pdf}
  \caption{Pairwise conditional reductions with the comparator named explicitly.}
\end{figure}
\FloatBarrier

\section{Results below the 9,000-m\textsuperscript{2} gate}
\eyebrow{Five distinct achieved humidity targets retain conditional savings}

Twelve converged eight-segment engineering-screen cases lie at or below the
area gate. Selecting the lowest known total power for each achieved humidity
ratio gives the following frontier.

\begin{center}\small
\begin{tabular}{@{}lrrrrr@{}}
\toprule
\textbf{Core outlet} & \textbf{Voltage} & \textbf{Area} & \textbf{Known total} & \textbf{Saving vs ERV+AC} & \textbf{Break-even} \\
 & [V] & [m\textsuperscript{2}] & [kW] & [kW / \%] & [Wh/kg] \\
\midrule
12.000~g/kg & 0.10 & 6,519 & 490.475 & 53.265 / 9.80 & 25.379 \\
11.000~g/kg & 0.10 & 8,669 & 446.484 & 97.257 / 17.89 & 35.272 \\
10.000~g/kg & 0.50 & 7,066 & 429.256 & 114.485 / 21.06 & 43.779 \\
9.000~g/kg & 1.00 & 5,758 & 422.938 & 120.802 / 22.22 & 51.172 \\
O30 (8.886~g/kg) & 1.00 & 5,899 & 418.400 & 125.340 / 23.05 & 51.955 \\
\bottomrule
\end{tabular}
\end{center}

\begin{figure}[htbp]
  \centering
  \includegraphics[width=0.99\textwidth]{../../outputs/integrated_erv_eo_under_9000_v1/under_9000_power_and_saving.pdf}
  \caption{Best integrated ERV+EO point at each achieved humidity ratio below 9,000~m\textsuperscript{2}.}
\end{figure}
\FloatBarrier

Deeper drying does not require monotonically increasing area because the
selected voltage changes. Increasing voltage raises flux and can reduce area,
but it also increases active-core power. The selected frontier is therefore a
joint voltage--area--energy result, not an area-only trend.

The 8.000-g/kg target needs at least approximately 12,735~m\textsuperscript{2}
among the valid eight-segment cases, while the full supply-humidity target
needs at least 17,253~m\textsuperscript{2}. Neither is accepted under the
current gate.

\section{Break-even and system-loss budget}
\eyebrow{The omitted-load allowance is finite and architecture-specific}

\begin{figure}[htbp]
  \centering
  \includegraphics[width=0.99\textwidth]{../../outputs/integrated_erv_eo_under_9000_v1/under_9000_area_and_break_even.pdf}
  \caption{Conditional saving and calculated active-core break-even for the selected under-gate cases. Orange is known active-core SEC; hatching is the remaining allowance.}
\end{figure}
\FloatBarrier

For O30, any combination of omitted receiver, heat-rejection, conversion,
electrode, manifold and incremental fan loads above 125.340~kW removes the
incremental advantage over passive ERV+AC. This power ceiling is equivalent to
the 28.709-Wh/kg gap between known active-core SEC and calculated break-even.
The limit must be treated as a design budget, not as a prediction that the
missing subsystems will fit inside it.

\begin{stopbox}
\textbf{No net-saving claim yet.} Every positive value in this section is a
known-power conditional window. The full-coupled-system gate is still open for
all selected under-9,000-m\textsuperscript{2} cases.
\end{stopbox}

\section{Model closure and membrane interpretation}
\eyebrow{The remaining uncertainty is empirical and subsystem-specific}

The model phase has resolved several interpretation errors:
\begin{itemize}
  \item a negative passive surface activity direction is not a thermodynamic
  rejection of electrically driven uphill EO transport;
  \item current density and membrane area are solved coordinates and must not
  be used as undocumented universal acceptance gates;
  \item water diffusivity alone does not select a membrane; the relevant
  quantity is a coupled bundle including sorption derivative, chemical water
  mobility/permeance, thickness, electro-osmotic drag, wet conductivity and
  compressed-stack resistance;
  \item a symmetric quasi-static $+J/0/-J$ replay produces no useful mean EO
  pumping beyond the passive state in the present constant-drag model;
  asymmetry, rectification or memory would require measured physics not yet in
  the model.
\end{itemize}

The direct active-core cases close the signed thermodynamic audit and the
selected under-gate cases pass the current transport, robust-physics and
engineering screens. They do not close receiver operation, conjugate thermal
feedback, real manifolds, leakage, durability or product-scale current
collection. A hypothetical joint material vector is therefore an experiment
target, not a procurement specification.

\begin{gobox}
\textbf{Phase decision: GO for evidence generation.} Stop unfocused parameter
sweeps. Acquire same-sample material data and run a matched EO-off/EO-on
integrated-ERV coupon experiment. Replay the model only after measured
properties and subsystem power replace the open assumptions.
\end{gobox}

\section{Vapour-cell experiment}
\eyebrow{Centimetre-scale controlled flow, not a shallow closed steady cell}

The recommended first fixture uses a 50$\times$50-mm active aperture
(25~cm\textsuperscript{2}), an approximately 80$\times$80-mm membrane coupon
and 110$\times$110-mm outer plates. Use interchangeable symmetric 1/2/4-mm
gaps, with 2~mm nominal. At 1~m/s the nominal flow is 6.0 actual L/min per
side. Use low-flow acquisition separately when humidity-sensor detectability
requires it. Start with 8--10-mm insulating flow frames plus 12-mm metal
backing/load plates; a load-bearing PEEK plate should start near 16~mm and be
accepted only after loaded flatness and gap mapping.

\begin{figure}[htbp]
  \centering
  \includegraphics[width=0.78\textwidth]{../../outputs/eo_vapor_cell_sizing_v1/cell_dimension_tradeoffs.png}
  \caption{Coupon-cell dimension trade-offs supporting the 25-cm\textsuperscript{2}, 2-mm-gap starting fixture.}
\end{figure}
\FloatBarrier

The main experimental controls are a split common inlet, ABBA polarity
reversal, two-sided water balance, foil blanks, pressure control, gap/flow
repeats, reference cross-switching and multiple aperture/perimeter inserts.
Static temperature or RH mismatch can dominate a one-direction estimate;
polarity-odd analysis rejects the static component but not nonlinear or
polarity-correlated drift. Closed chambers are transient inventory experiments
and must not be described as fixed-boundary steady-state measurements.

\section{Final decision and handoff}
\eyebrow{What is ready, what is not, and what happens next}

\begin{center}\small
\begin{tabularx}{\textwidth}{@{}p{54mm}p{22mm}Y@{}}
\toprule
\textbf{Decision} & \textbf{Status} & \textbf{Basis} \\
\midrule
Use passive ERV as primary control & \textcolor{Green}{\textbf{GO}} & Matched 543.741-kW equal-service reference \\
Test EO integrated inside ERV to O30 & \textcolor{Green}{\textbf{GO}} & 5,899~m\textsuperscript{2}; 125.340-kW conditional window \\
Use EO for O50 in this duty & \textcolor{Slate}{\textbf{NO DUTY}} & Passive ERV already passes O50 \\
Claim full EO dehumidification below 9,000~m\textsuperscript{2} & \textcolor{Red}{\textbf{NO-GO}} & Minimum valid screened full-duty area is about 17,253~m\textsuperscript{2} \\
Claim measured net energy saving & \textcolor{Orange}{\textbf{NOT YET}} & Receiver, heat rejection and hardware ledger remain open \\
Select a commercial membrane/module & \textcolor{Orange}{\textbf{NOT YET}} & Same-sample coupled properties and durability are unmeasured \\
Build controlled-flow coupon fixture & \textcolor{Green}{\textbf{GO}} & Geometry, controls and error mechanisms are specified \\
\bottomrule
\end{tabularx}
\end{center}

The highest-value next actions are:
\begin{enumerate}
  \item measure the same membrane's sorption, mobility/permeance, drag, wet
  conductivity, compressed ASR, thickness and mechanics;
  \item run matched EO-off/EO-on ABBA tests at controlled film conditions;
  \item close receiver, thermal, conversion and hydraulic power inside the
  125.340-kW O30 allowance;
  \item replay the integrated ERV model with measured data and decide whether
  the O30 window remains positive.
\end{enumerate}

\begin{mainbox}
\textbf{Final project conclusion.} The project has identified a credible
conditional opportunity and a concrete experiment capable of falsifying it.
The correct claim is not that EO has already saved energy, but that EO
integrated in the ERV can retain a 23.05\% known-power window at O30 within the
9,000-m\textsuperscript{2} gate, provided all omitted subsystem loads remain
below the calculated allowance. This is ready for technical handoff and
experimental continuation.
\end{mainbox}

\section{Controlled evidence basis}
\begin{itemize}
  \item \path{outputs/direct_active_eo_erv_study_v1/study/} --- integrated active-core case register and audits;
  \item \path{outputs/integrated_erv_eo_opportunity_plots_v2/} --- matched architecture, pairwise saving and full-duty screens;
  \item \path{outputs/integrated_erv_eo_under_9000_v1/} --- selected under-gate cases and calculated break-even;
  \item \path{outputs/eo_vapor_cell_sizing_v1/} and \path{outputs/eo_vapor_boundary_analysis_v1/} --- fixture sizing and boundary-condition analysis;
  \item \path{reports/weekly_findings_2026-08-24_to_2026-08-30/} --- Revision~4 weekly evidence chronology;
  \item \path{reports/eo_vapor_boundary_analysis_v2_en/} --- canonical English vapour-cell report.
\end{itemize}

\vfill
{\sans\scriptsize\color{Slate}End of controlled Version~9 report. All
positive energy differences are conditional model windows and do not
constitute material, module, product or commercial validation.}
\end{document}
